\subsection{非处处有定义的合成律的态射概念}

定义 1. 设 E、E' 是两个集合，各带一个记作 \((x,y)\mapsto x*y\) 的内合成律，它不必处处有定义。映射 \(f:E\to E'\) 称为态射，如果对 E 中所有使 \(f(x)*f(y)\) 有定义的 x、y，合成 \(x*y\) 也有定义，且：

\[
f (x * y) = f (x) * f (y).\tag{1}
\]

更简短地说，公式 (1) 必须在右端有意义的每一次都成立。

态射的概念不同于表示的概念（《代数》

第一章 § 1 第 1 节），在那里要求只要左端有意义，等式 (1) 就成立。当然，对于处处有定义的合成律，这两个概念是一致的。

定义 2. 设 E、E' 是两个集合，各带一族内合成律 \((x,y)\mapsto x*\alpha y,\alpha\in\mathbf{I}\)。映射 \(f:E\to E'\) 称为态射，如果它对每个合成律 \((x,y)\mapsto x*\alpha y\) 都是态射。

与表示一样，态射也满足《集合论》第四章 § 2 的公理 \((\mathrm{MO}_{\mathrm{I}})\)、\((\mathrm{MO}_{\mathrm{II}})\)、\((\mathrm{MO}_{\mathrm{III}})\)。若 \(f: E \to E'\) 是态射，则 \(f(E)\) 是 \(R'\) 的稳定子集。

